% En este archivo, incluímos cuestiones de formato, diseño, librerías, etc.

% Traduce automáticamente títulos ("Tabla", "Capítulo", etc.)
\usepackage[spanish,es-tabla]{babel}  

% Codificación UTF-8 (para tildes y ñ directamente)
\usepackage[utf8]{inputenc}    

% Control del interlineado
\usepackage{setspace}

% Sangría en el primer párrafo tras cada título
\usepackage{indentfirst} 

% Evita que el texto se salga de los márgenes
\sloppy

% Interlineado 1.5
\onehalfspacing 

% Espacio vertical entre párrafos
\setlength{\parskip}{10pt} 

% Para definir los márgenes
\usepackage{geometry}
\newgeometry{bottom=2cm, top=2cm, left=3cm, right=2cm}

% Definir el tipo de letra similar a Times New Roman
\usepackage[T1]{fontenc}
\usepackage{times}

% Más control en tablas (alineación, columnas especiales)
\usepackage{array}        

% Colores en texto (se carga abajo con la opción table; cargarlo dos veces causaba
% "Option clash for package xcolor")

% Colores en tablas (fondo de filas/columnas)
\usepackage[table]{xcolor}          

% Personalización de listas (itemize, enumerate)
\usepackage{enumitem}                  

% Insertar imágenes
\usepackage{graphicx}   

% Controlar posición de figuras/tablas con [H]
\usepackage{float}  

% Ecuaciones avanzadas
\usepackage{amsmath}                   
% Mostrar URLs con saltos de línea automáticos
\usepackage{url}       

% Comentar bloques grandes de texto
\usepackage{comment}                   

% Personalizar encabezados y pies
\usepackage{fancyhdr}      

% Activar estilo fancy
\pagestyle{fancy}                      
% Limpiar encabezado y pie
\fancyhf{}                         

% Encabezado: capítulo (en report/book)
\fancyhead[LE, RO]{\small\leftmark} 

% Encabezado: número de página
\fancyhead[RE,LO]{\small\thepage}      

% Ajustar altura del encabezado para evitar warnings
\setlength{\headheight}{14pt}        

% Permite dibujar líneas de cualquier tamaño y colocarlas exactamente donde sea
\usepackage{tikz}

% Esta librería se usó para aumentar un poco el margen izquierdo de la portada
\usepackage{changepage}

% Comentar varias lineas
\usepackage{comment}